\chapter{编写 Markdown 混排内容}\label{ch:markdown}
\section{正文、标题与列表}
空行分隔段落。标题使用一至六个井号，井号后留一个空格。列表、引用和链接可按下例编写；示例含中文和标点，需先导入相应轨迹。
\par\noindent\begin{minipage}{\linewidth}
\begin{Source}
# 测量记录

这是第一段。下一段前留一个空行。

## 操作步骤

1. 检查接线
2. 记录读数

- 第一次测量
- 第二次测量

> 注意单位一致。

[查看说明](https://example.com)
\end{Source}
\end{minipage}\par\medskip

链接只排版可见文字，不把目标网页带入页面。强调、加粗和行内代码保留语义，但手写资产不会因此自动变成系统字体的粗体或斜体。任意 HTML 与 CSS 不是本编辑器的排版接口；需要控制页面位置时使用模板。

\section{任务列表}
\par\noindent\begin{minipage}{\linewidth}
\begin{Source}
- [ ] 12.50
- [x] 13.25
\end{Source}
\end{minipage}\par\medskip
方框和勾线由可书写折线绘制。点击预览中的方框会更新源码，可使用正文撤销恢复。列表文字仍需要字形资源；方框不需要另外采集。只有当前预览完成后才可操作，过期页面不会修改新源码。

\Needspace{8\baselineskip}
\section{脚注}
\entry{设置 → 元素 → 脚注设置…}
\par\noindent\begin{minipage}{\linewidth}
\begin{Source}
12.50[^1]

[^1]: 0.05
\end{Source}
\end{minipage}\par\medskip
脚注按首次引用顺序编号，放在页底；同一标签多次引用共享注释。长注释可续页。设置可调整标记与注释字号、行高、与正文的距离、分隔线以及页底预留比例。

定义放在顶层。脚注可以包含已支持的代码、公式、表格和图形；复杂结构仍遵循各自的分页规则。不要在定义内递归引用另一个脚注。缺失定义或不可容纳的内容会显示诊断，先修正再导出。

\section{代码与字面内容}
\entry{设置 → 元素 → 代码设置…}
使用单反引号包住行内代码，或使用三反引号围成代码块。代码区不解释数学定界符或普通 Markdown 标记。例如：
\par\noindent\begin{minipage}{\linewidth}
\begin{Source}
```text
12.50
  13.25
```
\end{Source}
\end{minipage}\par\medskip
代码块保留空行、空格和制表位。默认字高 4 mm、字格宽 5 mm、最小行高 6 mm、制表位 4 列。默认超宽行报错；可在设置中明确允许换行，续行会带标记。程序代码中的英文字母和符号也必须有真实字形资源，不自动使用等宽印刷字体作为机器轨迹。

\section{图片与单线图}
\entry{设置 → 元素 → 图片设置…}
先在源码中独占一段写图片标记，再在图片设置中选择这一项并导入本地文件。
图片设置的应用范围可选择文档默认或任意一张图片。默认设置用于沿用默认的图片及后续新增图片，已有独立设置不受影响。选择具体图片后可设置个性化参数，或勾选恢复沿用文档默认；设置及资源随项目保存，取消不写入项目。没有图片时也可以预先配置默认尺寸。
\par\noindent\begin{minipage}{\linewidth}
\begin{Source}
![图一 测量示意](assets/chart.png)
\end{Source}
\end{minipage}\par\medskip
方括号文字作为图题；无文字则不显示图题。图题不自动编号。0.1.24 起默认高 60 mm，宽度按原图比例计算，图题间距 2 mm、居中；默认页不再提供宽度输入。单张图片可在右侧独立选择按高度或按宽度等比设置；旧文档尺寸保持原值。图片与图题整体分页，不支持正文环绕或表格单元格内图片。

\par\noindent\begin{tabularx}{\linewidth}{@{}P{31mm}Y Y@{}}
\toprule
导入类型 & 使用方式 & 可用输出\\\midrule
PNG / JPEG & 作为完整图片，与手写图题一起预览。 & 参考 SVG/PNG；机器 SVG 导出前提示跳过图片，保留占位与图题。\\
受控单线 SVG & 导入作者确认的中心线路径，设置纸面尺寸和曲线容差。 & 校验通过后可与正文一起导出机器路径。\\
模板底图 & 在模板设计器中导入，帮助定位预印纸。 & 仅供参考，不作为要书写的正文图像。\\
\bottomrule
\end{tabularx}

单线 SVG 需要无填充、明确黑色描边和有效宽高；不接受文字元素、外部链接、脚本、滤镜、虚线和任意字体轮廓。曲线按最终纸面误差转换为折线，默认容差 0.05 mm。调小容差可增加细节，也会增加路径点。导入失败时应按诊断修改原图，不把更改扩展名当作格式转换。

\section{Obsidian 静态扩展}
\entry{设置 → 正文 → Markdown 格式 → Obsidian 静态扩展}
\par\noindent\begin{minipage}{\linewidth}
\begin{Source}
[[资料/测量|12.50]]

> [!note]- 测量记录
> 13.25

![[附件/原记录|原记录说明]]
\end{Source}
\end{minipage}\par\medskip
双链只显示标签；无别名时显示目标文字。Callout 按引用样式展开，折叠标记不会把纸面内容隐藏。带显式标签的嵌入只显示标签，并提醒目标内容未展开；没有标签的嵌入会报错。该模式不读取 Obsidian 库，也不执行插件、CSS、动态查询或递归嵌入。

管道表格中的别名分隔符应转义，例如 \code{[[资料/测量\|12.50]]}。需要真正放入图片时，使用普通 Markdown 图片和显式资源导入流程。
